% Options for packages loaded elsewhere
\PassOptionsToPackage{unicode}{hyperref}
\PassOptionsToPackage{hyphens}{url}
\documentclass[
]{article}
\usepackage{xcolor}
\usepackage{amsmath,amssymb}
\setcounter{secnumdepth}{-\maxdimen} % remove section numbering
\usepackage{iftex}
\ifPDFTeX
  \usepackage[T1]{fontenc}
  \usepackage[utf8]{inputenc}
  \usepackage{textcomp} % provide euro and other symbols
\else % if luatex or xetex
  \usepackage{unicode-math} % this also loads fontspec
  \defaultfontfeatures{Scale=MatchLowercase}
  \defaultfontfeatures[\rmfamily]{Ligatures=TeX,Scale=1}
\fi
\usepackage{lmodern}
\ifPDFTeX\else
  % xetex/luatex font selection
\fi
% Use upquote if available, for straight quotes in verbatim environments
\IfFileExists{upquote.sty}{\usepackage{upquote}}{}
\IfFileExists{microtype.sty}{% use microtype if available
  \usepackage[]{microtype}
  \UseMicrotypeSet[protrusion]{basicmath} % disable protrusion for tt fonts
}{}
\makeatletter
\@ifundefined{KOMAClassName}{% if non-KOMA class
  \IfFileExists{parskip.sty}{%
    \usepackage{parskip}
  }{% else
    \setlength{\parindent}{0pt}
    \setlength{\parskip}{6pt plus 2pt minus 1pt}}
}{% if KOMA class
  \KOMAoptions{parskip=half}}
\makeatother
\usepackage{longtable,booktabs,array}
\newcounter{none} % for unnumbered tables
\usepackage{calc} % for calculating minipage widths
% Correct order of tables after \paragraph or \subparagraph
\usepackage{etoolbox}
\makeatletter
\patchcmd\longtable{\par}{\if@noskipsec\mbox{}\fi\par}{}{}
\makeatother
% Allow footnotes in longtable head/foot
\IfFileExists{footnotehyper.sty}{\usepackage{footnotehyper}}{\usepackage{footnote}}
\makesavenoteenv{longtable}
\usepackage{graphicx}
\makeatletter
\newsavebox\pandoc@box
\newcommand*\pandocbounded[1]{% scales image to fit in text height/width
  \sbox\pandoc@box{#1}%
  \Gscale@div\@tempa{\textheight}{\dimexpr\ht\pandoc@box+\dp\pandoc@box\relax}%
  \Gscale@div\@tempb{\linewidth}{\wd\pandoc@box}%
  \ifdim\@tempb\p@<\@tempa\p@\let\@tempa\@tempb\fi% select the smaller of both
  \ifdim\@tempa\p@<\p@\scalebox{\@tempa}{\usebox\pandoc@box}%
  \else\usebox{\pandoc@box}%
  \fi%
}
% Set default figure placement to htbp
\def\fps@figure{htbp}
\makeatother
\setlength{\emergencystretch}{3em} % prevent overfull lines
\providecommand{\tightlist}{%
  \setlength{\itemsep}{0pt}\setlength{\parskip}{0pt}}
\usepackage{bookmark}
\IfFileExists{xurl.sty}{\usepackage{xurl}}{} % add URL line breaks if available
\urlstyle{same}
\hypersetup{
  hidelinks,
  pdfcreator={LaTeX via pandoc}}

\author{}
\date{}

\begin{document}

\textbf{KYAMBOGO}\includegraphics[width=1.54722in,height=1.16944in]{./media/image1.png}
\textbf{UNIVERSITY}

\textbf{BACHELOR OF ENGINEERING IN MECHATRONICS AND BIOMEDICAL
ENGINEERING.}

\textbf{TEMB 1202: COMPUITING II}

\textbf{ASSIGNMENT: EMPLOYEE PAYROLL SYSTEM}

\textbf{DATE: 27\textsuperscript{TH} MAY, 2026}

\textbf{GROUP D EVENING}

{\def\LTcaptype{none} % do not increment counter
\begin{longtable}[]{@{}
  >{\raggedright\arraybackslash}p{(\linewidth - 2\tabcolsep) * \real{0.4952}}
  >{\raggedright\arraybackslash}p{(\linewidth - 2\tabcolsep) * \real{0.5048}}@{}}
\toprule\noalign{}
\begin{minipage}[b]{\linewidth}\centering
\textbf{NAME}
\end{minipage} & \begin{minipage}[b]{\linewidth}\centering
\textbf{REGISTRATION NUMBER}
\end{minipage} \\
\midrule\noalign{}
\endhead
\bottomrule\noalign{}
\endlastfoot
KATONGOLE SHADDAI KISA & 25/U/BIE/01380/PE \\
ZALWANGO CHRISTINE & 25/U/BIE/01426/PE \\
MUTYABA DERICK & 25/U/BIE/01399/PE \\
NAKANWAGI MARIAM & 25/U/BIE/01403/PE \\
BUJINGO DAVID KABANDA & 25/U/BIE/01364/PE \\
RWOTOMIA PIUS EMMANUEL & 25/U/BIE/01419/PE \\
\end{longtable}
}

\textbf{INTRODUCTION}

This report documents the design and implementation of an Employee
Payroll System written in the C programming language. The system is a
fully menu-driven console application that enables payroll officers to
manage employee records, compute salaries with overtime, apply statutory
deductions, and generate comprehensive payroll reports. It demonstrates
practical application of C programming concepts including structures,
arrays, functions, pointers, loops, and conditional logic. The program
runs continuously until the user selects Exit, and requires login
authentication before any payroll operations can be performed. All
inputs are validated against defined business rules to prevent incorrect
salary calculations

\textbf{PROBLEM STATEMENT}

An organization needs a digital payroll system to replace manual salary
processing. The system must handle the complete payroll lifecycle from
registering employees to computing gross pay, applying deductions,
calculating net pay, and producing summary reports while enforcing the
company\textquotesingle s payroll rules.

\textbf{Required Operations}

\begin{itemize}
\item
  Register new employees with name, department, hourly rate, and hours
  worked.
\item
  Search existing employee records by name or employee ID.
\item
  Update employee details including rate, hours worked, and department.
\item
  Delete employee records for staff who have left the organization.
\item
  Calculate salaries, gross pay, overtime, deductions, and net pay.
\item
  Display a complete payroll summary report identifying the highest-paid
  employee.
\end{itemize}

\textbf{Business rules}

{\def\LTcaptype{none} % do not increment counter
\begin{longtable}[]{@{}
  >{\raggedright\arraybackslash}p{(\linewidth - 2\tabcolsep) * \real{0.5000}}
  >{\raggedright\arraybackslash}p{(\linewidth - 2\tabcolsep) * \real{0.5000}}@{}}
\toprule\noalign{}
\begin{minipage}[b]{\linewidth}\raggedright
Normal hours limit
\end{minipage} & \begin{minipage}[b]{\linewidth}\raggedright
Normal working hours are capped at 160 hours/month.
\end{minipage} \\
\midrule\noalign{}
\endhead
\bottomrule\noalign{}
\endlastfoot
Overtime pays & Hours above 160 are paid at 1.5* the base hourly
rate. \\
Excessive hours warning & Hours above 300 trigger a visible warning
message. \\
Deductions on Gross pay & Tax 30\%, NSSF 5\%, Health insurance 2\% of
gross pay. \\
Negative input rejection & Negative hours or rates are rejected with an
error. \\
Highest-paid employee & The report identifies the employee with highest
net pay. \\
\end{longtable}
}

\textbf{SYSTEM FEATURES AND FUNCTIONALITY}

Authentication On launch the system presents a login prompt. After three
consecutive failed attempts access is denied and the program exits,
preventing unauthorized access.

\textbf{Main menu}

{\def\LTcaptype{none} % do not increment counter
\begin{longtable}[]{@{}
  >{\raggedright\arraybackslash}p{(\linewidth - 2\tabcolsep) * \real{0.4911}}
  >{\raggedright\arraybackslash}p{(\linewidth - 2\tabcolsep) * \real{0.4911}}@{}}
\toprule\noalign{}
\begin{minipage}[b]{\linewidth}\raggedright
Feature
\end{minipage} & \begin{minipage}[b]{\linewidth}\raggedright
Description
\end{minipage} \\
\midrule\noalign{}
\endhead
\bottomrule\noalign{}
\endlastfoot
Registration employee & Add a new employee record \\
Search employee & Find by name or ID \\
Update employee & Modify rate, hours, or department \\
Delete employee & Remove a record permanently \\
Calculate all salaries & Refresh payroll for all employees \\
Display payroll report & Full summary + highest-paid employee \\
Exit & Logout and close the system \\
\end{longtable}
}

\textbf{Salary Calculation Steps}

{\def\LTcaptype{none} % do not increment counter
\begin{longtable}[]{@{}
  >{\raggedright\arraybackslash}p{(\linewidth - 4\tabcolsep) * \real{0.3202}}
  >{\raggedright\arraybackslash}p{(\linewidth - 4\tabcolsep) * \real{0.3203}}
  >{\raggedright\arraybackslash}p{(\linewidth - 4\tabcolsep) * \real{0.3203}}@{}}
\toprule\noalign{}
\begin{minipage}[b]{\linewidth}\raggedright
Step
\end{minipage} & \begin{minipage}[b]{\linewidth}\raggedright
Item
\end{minipage} & \begin{minipage}[b]{\linewidth}\raggedright
Formula
\end{minipage} \\
\midrule\noalign{}
\endhead
\bottomrule\noalign{}
\endlastfoot
Step 1 & Normal pay & Hours \textless=160 - Hours × Rate \\
Step 2 & Overtime pays & (Hours - 160) × Rate × 1.5 \\
Step 3 & Gross pay & Normal pay + Overtime pay \\
Step 4 & Tax & Gross × 30\% \\
Step 5 & NSSF & Gross × 5\% \\
Step 6 & Health ins & Gross × 2\% \\
Step 7 & Total deductions & Tax + NSSF + Health insurance \\
Step 8 & Net pay & Gross pay -- Total deductions \\
\end{longtable}
}

\textbf{PROGRAM DESIGN AND STRUCTURE}

\textbf{Data Structure}

The central data structure is a C \textbf{struct} called
\textbf{Employee}. An array of up to MAX\_EMPLOYEES (100) of this struct
acts as the in-memory database.

typedef struct \{

int empID;

char name{[}60{]};

char department{[}40{]};

float hourlyRate;

float hoursWorked;

float normalPay;

float overtimePay;

float grossPay;

float taxAmount;

float nssfAmount;

float healthAmount;

float totalDeductions;

float netPay;

int overtimeFlag;

int excessiveFlag;

\} Employee;

\textbf{Constants}

\#define MAX\_EMPLOYEES 100

\#define NORMAL\_HOURS 160.0f

\#define EXCESSIVE\_HOURS 300.0f

\#define OT\_MULTIPLIER 1.5f

\#define TAX\_RATE 0.30f

\#define NSSF\_RATE 0.05f

\#define HEALTH\_RATE 0.02f

\#define ID\_BASE 100

\textbf{IMPORTANT FUNCTIONS USED}

login()

Signature: void login(void)

Handles authentication. Allows 3 attempts then exits. First function
called in main().

registerEmployee()

Signature: void registerEmployee(void)

Collects employee data, validates negatives, auto-assigns ID, calls
calculateSalary().

searchEmployee()

Signature: void searchEmployee(void)

Finds by case-insensitive name (partial match) or exact ID. Displays
full record.

updateEmployee()

Signature: void updateEmployee(void)

Updates rate, hours, department, or all fields. Recalculates salary
after change.

deleteEmployee()

Signature: void deleteEmployee(void)

Confirms deletion, removes record by array shifting, decrements count.

calculateSalary()

Signature: void calculateSalary(Employee *emp)

Core computation: normal pay, overtime, gross pay, all deductions, net
pay, flags.

calculateAllSalaries()

Signature: void calculateAllSalaries(void)

Loops through all employees and calls calculateSalary() on each record.

displayPayrollReport()

Signature: void displayPayrollReport(void)

Prints formatted table, totals, highest-paid employee, and deduction
rates used.

displayEmployee()

Signature: void displayEmployee(Employee *emp)

Prints a single employee\textquotesingle s complete payroll breakdown
neatly.

findByID()

Signature: int findByID(int id)

Linear search returning array index for matching empID, or -1 if not
found.

findByName()

Signature: int findByName(const char *query)

Case-insensitive partial string match using strstr(). Returns first
matching index.

clearBuffer()

Signature: void clearBuffer(void)

Flushes stdin after scanf() to prevent stale newline issues.

toLowerCase() Signature: void toLowerCase(char *str, char *out)

Converts string to lowercase for case-insensitive comparisons.

printDivider()

Signature: void printDivider(char ch, int len)

Prints a horizontal line of a given character --- used for consistent UI
formatting

\textbf{DESIGN DECISIONS AND JUSTIFICATIONS}

Storing All Computed Fields in the Struct

grossPay, deductions, netPay, overtimeFlag, and excessiveFlag are stored
in the Employee struct rather than recalculated on demand. This ensures
every display function reads consistent values and avoids duplicating
the calculation logic.

\#define Constants for All Rates

Using named constants (TAX\_RATE, OT\_MULTIPLIER, etc.) rather than
magic numbers makes the code self-documenting. Changing a statutory rate
requires editing one line only, which is essential for payroll software
where rates change annually.

calculateSalary() Called at Registration and After Updates

Isolating the calculation in its own function and calling it both on
registration and after every update guarantees that stored figures are
always current. Without this, updating an hourly rate would leave stale
net pay values.

Employee IDs Starting at 1001 Starting IDs at 1001 produces
professional-looking employee numbers clearly distinct from array
indices, reducing off-by-one confusion.

fgets for Name Input

fgets() reads full lines including spaces, so "John Doe" is captured
correctly instead of being truncated at the space by scanf("\%s").

Global Array Instead of Linked List

An array is simpler to implement using the C concepts covered in this
course, provides O(1) indexed access, and is adequate for up to 100
employees. A linked list would be preferable for very large, dynamic
datasets

\textbf{CHALLENGES ENCOUNTERED}

Reading Names with Spaces

scanf("\%s") stops at spaces. We switched to fgets() and stripped the
trailing newline with strcspn() to correctly capture full names like
"John Doe".

Stale Salary After Updates

Early versions did not recalculate salary after an update. We fixed this
by always calling calculateSalary() at the end of updateEmployee().

Input Buffer Contamination

Mixing scanf() for numbers and fgets() for strings caused empty reads. A
clearBuffer() helper that consumes remaining characters resolved this.

Case-Insensitive Search

The standard strcmp() is case-sensitive. We implemented toLowerCase()
and used strstr() to support partial, case-insensitive name searches.

Safe Report When List is Empty

Accessing employees{[}0{]} when empCount is 0 causes undefined
behaviour. We added a guard (if empCount == 0) before entering the
report loop.

GitHub Merge Conflicts

Multiple members editing the same struct definition caused conflicts. We
resolved this by agreeing on the final struct layout first and
committing it before branching for individual functions

\textbf{CONCLUSION}

The Employee Payroll System was successfully designed, implemented, and
tested as a fully functional menu-driven C program satisfying all
requirements of the assignment. All business rules were enforced: normal
hours are capped at 160/month, overtime is paid at 1.5× the base rate,
hours above 300 trigger a visible warning, deductions (tax 30\%, NSSF
5\%, health 2\%) are computed on gross pay, negative inputs are
rejected, and the payroll report correctly identifies the highest-paid
employee. This project provided valuable hands-on experience translating
real-world business requirements into working software, reinforced core
C programming skills, and built collaborative development habits through
GitHub version control.

\end{document}
